\section*{Introducción}
\addcontentsline{toc}{section}{Introducción}

En la actualidad, el monitoreo ambiental es una tarea fundamental para comprender y anticipar cambios en el clima, la calidad del aire y la variabilidad de condiciones atmosféricas. Las estaciones meteorológicas permiten registrar variables como temperatura, humedad, presión atmosférica, velocidad y dirección del viento, entre otras, con el fin de apoyar la toma de decisiones en ámbitos agrícolas, académicos, de investigación y vigilancia ambiental.

Con el avance de la electrónica y el Internet de las Cosas, es posible diseñar sistemas más accesibles y económicos para la captación y publicación de datos en tiempo real. En este contexto, el uso de microcontroladores como Arduino facilita la adquisición de información desde sensores, mientras que plataformas web permiten visualizar, consultar y analizar la información de manera centralizada.

Este proyecto propone el diseño e implementación de un sistema web para el monitoreo de estaciones meteorológicas basadas en Arduino. La solución integra la recolección automática de datos, su almacenamiento en una base de datos y la presentación de información mediante una interfaz web con acceso a indicadores y gráficos. De esta manera, se busca ofrecer una herramienta útil para la observación continua del entorno, la comparación histórica de datos y la generación de información de apoyo para usuarios finales e investigadores.

El desarrollo del sistema contempla tanto la parte de captura y procesamiento de datos como la parte de interacción con los usuarios, con el objetivo de facilitar la consulta de información climática en cualquier momento y desde distintos dispositivos. Además, la implementación de una arquitectura basada en tecnologías web modernas permite mejorar la escalabilidad, la mantenibilidad y la disponibilidad del servicio para su uso en entornos reales.

En conjunto, este trabajo representa una propuesta tecnológica orientada a la automatización del monitoreo meteorológico, promoviendo la integración de sensores electrónicos, almacenamiento de información y visualización en línea como elementos clave para la análisis y gestión del entorno.

\subsection*{Breve historia de la Meteorología}
Antes de todos los tiempos, cuando no existian instrumentos de medición, los seres humanos dependían de la observación directa de la naturaleza para predecir el clima. Los antiguos griegos y romanos ya tenían nociones básicas sobre los fenómenos meteorológicos, aunque carecían de herramientas precisas para su estudio.

\subsubsection*{El aporte de Von Newman}
Cuando se creo la primera computadora en 1945, Von Newman propuso un modelo matemático para predecir el clima. Este modelo se basaba en ecuaciones diferenciales que describían la dinámica de la atmósfera y los océanos. A partir de este modelo, se desarrollaron los primeros sistemas de predicción numérica del tiempo, que permitieron realizar pronósticos más precisos y confiables.